\section{Dwa konkretne obroty generujące grupę wolną}
\subsection{Grupa obrotów i punkty stałe jej elementów}\label{subsec:so3}
\begin{definicja}
Macierz $Q\in M_3(\R)$ jest ortogonalna, jeżeli $Q^TQ=I$. Definiujemy $\mathrm{SO}(3)=\{Q\in M_3(\R):Q^TQ=I,\ \det Q=1\}$. Macierze te zachowują długości i orientację. Działają na $\R^3$ przez $x\mapsto Qx$ oraz, przez ograniczenie, na $S^2$.
\end{definicja}
Rzeczywiście $\|Qx\|^2=x^TQ^TQx=\|x\|^2$. Iloczyn dwóch takich macierzy i odwrotność $Q^{-1}=Q^T$ również spełniają warunki definicji. Mamy więc grupę.
\begin{lemat}[Oś nietrywialnego obrotu]\label{lem:os-obrotu}
Jeśli $Q\in\mathrm{SO}(3)$ i $Q\ne I$, to zbiór wektorów ustalanych przez $Q$ jest prostą przez $0$. W szczególności $Q$ ma dokładnie dwa punkty stałe na $S^2$.
\end{lemat}
\begin{proof}
Najpierw wykażemy istnienie niezerowego wektora stałego. Mamy
\[
 \det(Q-I)=\det Q\det(I-Q^{-1})=\det(I-Q^T)
 =\det(I-Q)=-\det(Q-I).
\]
Ostatni znak minus wynika z nieparzystego wymiaru $3$. Zatem $\det(Q-I)=0$, więc układ $(Q-I)x=0$ ma rozwiązanie $x\ne0$.

Gdyby $Q$ ustalało dwa liniowo niezależne wektory, ustalałoby całą rozpiętą przez nie płaszczyznę. Wybierzmy w niej ortonormalną bazę $u,v$ i jednostkowy wektor $w$ prostopadły do tej płaszczyzny. Ortogonalność implikuje $Qw\perp u,v$ i $\|Qw\|=1$, zatem $Qw=\pm w$. W bazie $u,v,w$ wyznacznik wynosi odpowiednio $1$ albo $-1$. Warunek $\det Q=1$ wymusza $Qw=w$, a wtedy $Q=I$. To sprzeczność. Przestrzeń punktów stałych ma więc wymiar dokładnie jeden. Prosta przez $0$ przecina sferę w dwóch punktach antypodycznych.
\end{proof}
\subsection{Wybór generatorów}
Rozważmy
\begin{equation}\label{eq:5-1}
 A=\frac13\begin{pmatrix}3&0&0\\0&1&-2\sqrt2\\0&2\sqrt2&1\end{pmatrix},\qquad
 B=\frac13\begin{pmatrix}1&-2\sqrt2&0\\2\sqrt2&1&0\\0&0&3\end{pmatrix}.
\end{equation}
W obu macierzach blok $2\times2$ ma postać $\begin{psmallmatrix}c&-s\\s&c\end{psmallmatrix}$, gdzie $c=1/3$, $s=2\sqrt2/3$ i $c^2+s^2=1$. Jego kolumny są ortonormalne i wyznacznik wynosi $1$. Wobec tego $A,B\in\mathrm{SO}(3)$. Macierz $A$ obraca wokół osi $x$, a $B$ wokół osi $z$, o kąt o cosinusie $1/3$.

Z własności uniwersalnej istnieje homomorfizm $\rho:F_2\to\mathrm{SO}(3)$, $\rho(a)=A$, $\rho(b)=B$. Musimy udowodnić, że jest \emph{różnowartościowy}. To znacznie mocniejsze niż $AB\ne BA$.

\subsection{Rachunek modulo 3 i usunięcie pierwiastków}
Dla liczb całkowitych $m\equiv n\pmod3$ oznacza, że $m-n$ jest podzielne przez $3$. Występują tylko trzy reszty, które wygodnie zapisujemy jako $0,1,-1$. Mamy $2\equiv-1$. Redukcja modulo $3$ zachowuje sumy i iloczyny. Dotyczy to również iloczynów macierzy całkowitych, ponieważ ich wpisy powstają przez sumowanie iloczynów wpisów. Reszty tworzą ciało $\F_3$; potrzebujemy jedynie tego, że iloczyn dowolnej liczby czynników $1$ i $-1$ nie jest zerem. Nie będziemy dzielić przez $3$ po redukcji modulo $3$.

Dla litery $s\in\{a,\bar a,b,\bar b\}$ oznaczmy przez $R_s$ odpowiednio $A,A^{-1},B,B^{-1}$. Macierze $3R_s$ mają wpisy z pierwiastkiem $\sqrt2$. Aby dostać macierze całkowite, wprowadzamy pomocniczo $D_0=\operatorname{diag}(1,\sqrt2,1)$ oraz $N_s=3D_0^{-1}R_sD_0$. Bezpośrednie mnożenie daje
\begin{equation}\label{eq:5-2}
 N_a=\begin{pmatrix}3&0&0\\0&1&-2\\0&4&1\end{pmatrix},\qquad
 N_{\bar a}=\begin{pmatrix}3&0&0\\0&1&2\\0&-4&1\end{pmatrix},
\end{equation}
\begin{equation}\label{eq:5-3}
 N_b=\begin{pmatrix}1&-4&0\\2&1&0\\0&0&3\end{pmatrix},\qquad
 N_{\bar b}=\begin{pmatrix}1&4&0\\-2&1&0\\0&0&3\end{pmatrix}.
\end{equation}
Na przykład wpis $(2,3)$ w $D_0^{-1}(3A)D_0$ wynosi $(1/\sqrt2)(-2\sqrt2)\cdot1=-2$, a wpis $(3,2)$ wynosi $1\cdot2\sqrt2\cdot\sqrt2=4$. Macierze odwrotnych obrotów otrzymujemy z $A^{-1}=A^T$, $B^{-1}=B^T$, co wyjaśnia zmianę znaków.
\begin{uwaga}[Rola pomocniczej zmiany bazy]
Macierz $D_0$ nie jest obrotem i służy wyłącznie do obliczeń. Sprawdzamy za jej pomocą relacje między obrotami $A,B$. W konstrukcji kawałki sfery i kuli przemieszczamy jedynie obrotami i przesunięciami.
\end{uwaga}

\subsection{Cztery wektory i jedna tabela}
W rachunku modulo $3$ przyjmijmy wektory kolumnowe
\[
 w_a=\begin{pmatrix}0\\1\\1\end{pmatrix},\quad
 w_{\bar a}=\begin{pmatrix}0\\1\\-1\end{pmatrix},\quad
 w_b=\begin{pmatrix}1\\-1\\0\end{pmatrix},\quad
 w_{\bar b}=\begin{pmatrix}1\\1\\0\end{pmatrix}.
\]
Iloczyn $uv^T$ oznacza macierz, której wpis $(i,j)$ to $u_iv_j$. Z \eqref{eq:5-2}--\eqref{eq:5-3} otrzymujemy
\begin{equation}\label{eq:5-4}
 \overline{N_s}=w_sw_s^T,
\end{equation}
gdzie kreska oznacza redukcję wszystkich wpisów modulo $3$. Na przykład $\overline{N_a}=\begin{psmallmatrix}0&0&0\\0&1&1\\0&1&1\end{psmallmatrix}=w_aw_a^T$ i $\overline{N_b}=\begin{psmallmatrix}1&-1&0\\-1&1&0\\0&0&0\end{psmallmatrix}=w_bw_b^T$. Pozostałe dwie równości sprawdzamy identycznie.

Tabela iloczynów skalarnych $w_s^Tw_t$, już modulo $3$, ma postać
\begin{equation}\label{eq:5-5}
\begin{array}{c|rrrr}
 s\backslash t&a&\bar a&b&\bar b\\\hline
 a&-1&0&-1&1\\
 \bar a&0&-1&-1&1\\
 b&-1&-1&-1&0\\
 \bar b&1&1&0&-1
\end{array}
\end{equation}
Na przekątnej jest $1+1=2\equiv-1$. Dla pary $a,\bar a$ mamy $1-1=0$, a dla $b,\bar b$ także $1-1=0$. Przy dwóch różnych osiach tylko jedna współrzędna daje niezerowy wkład, równy $1$ albo $-1$. Otrzymujemy ważną równoważność:
\begin{equation}\label{eq:5-6}
 w_s^Tw_t=0\quad\Longleftrightarrow\quad t=s^{-1} \qquad(\text{w }\F_3).
\end{equation}
\begin{lemat}[Iloczyn pomocniczych macierzy słowa zredukowanego]\label{lem:liczniki}
Jeżeli $s_1\cdots s_n$ jest niepustym słowem zredukowanym, to $\overline{N_{s_1}\cdots N_{s_n}}$ nie jest macierzą zerową.
\end{lemat}
\begin{proof}
Dla wektorów $u,v,x,y$ zachodzi $(uv^T)(xy^T)=u(v^Tx)y^T$. Wpis $(i,j)$ po lewej stronie jest równy $\sum_k u_iv_kx_ky_j=u_i(\sum_kv_kx_k)y_j$, co dowodzi wzoru bez żadnej dodatkowej teorii. Stosując go kolejno do \eqref{eq:5-4}, dostajemy
\begin{equation}\label{eq:5-7}
 \overline{N_{s_1}\cdots N_{s_n}}=
 \left(\prod_{j=1}^{n-1}w_{s_j}^Tw_{s_{j+1}}\right)w_{s_1}w_{s_n}^T.
\end{equation}
Dla $n=1$ iloczyn skalarnych czynników jest pusty i z definicji równy $1$. Ponieważ słowo jest zredukowane, $s_{j+1}\ne s_j^{-1}$. Według \eqref{eq:5-6} każdy czynnik skalarny jest $1$ albo $-1$, więc ich iloczyn również jest niezerowy. Każdy z wektorów $w_s$ ma pewną współrzędną równą $1$. Macierz $w_{s_1}w_{s_n}^T$ ma zatem co najmniej jeden wpis równy $1$. Pomnożenie przez niezerowy skalar nie może wyzerować tego wpisu.
\end{proof}
\begin{twierdzenie}[Wolna podgrupa obrotów]\label{tw:wolne-obroty}
Homomorfizm $\rho:F_2\to\mathrm{SO}(3)$ określony przez $\rho(a)=A$, $\rho(b)=B$ jest różnowartościowy.
\end{twierdzenie}
\begin{proof}
Niech $w=s_1\cdots s_n$ będzie niepustym słowem zredukowanym. Z definicji $N_s$ i kolejnych skróceń $D_0D_0^{-1}$ wynika $N_{s_1}\cdots N_{s_n}=3^nD_0^{-1}\rho(w)D_0$. Gdyby $\rho(w)=I$, prawa strona byłaby równa $3^nI$. Ponieważ $n\ge1$, jej redukcja modulo $3$ byłaby zerowa, wbrew lematowi~\ref{lem:liczniki}. Zatem żadne niepuste słowo zredukowane nie przechodzi na $I$.

Jeśli $\rho(u)=\rho(v)$, to $\rho(v^{-1}u)=I$. Jedynym elementem mającym taki obraz jest $e$, więc $v^{-1}u=e$ i $u=v$. To dokładnie różnowartościowość.
\end{proof}
\begin{uwaga}[Czego rachunek modulo 3 nie dowodzi]
Nie zanurzyliśmy nieskończonej grupy w skończonej grupie macierzy modulo $3$. Macierze $\overline{N_s}$ są osobliwe. Ponadto $N_sN_{s^{-1}}=9I$, a nie $I$. Badaliśmy podzielność wpisów liczników iloczynu obrotów, a nie homomorfizm grupy obrotów uzyskany przez redukcję ułamka $1/3$.
\end{uwaga}

